\chapter{Objects and Messages}
\label{ch:objects}

\standfirst{Everything is an object and the only thing you can do to one is
send it a message. That sentence is not a slogan; it is a complete
specification of the computational model, and this chapter takes it
literally.}

\section{Everything}

There is no primitive type, no built-in operator, no value that is not an
object. \ct{3} is an instance of \ct{SmallInteger}. \ct{true} is the sole
instance of \ct{True}. \ct{nil} is the sole instance of \ct{UndefinedObject}.
Classes are objects; so are the methods inside them, the blocks you write, and
the stack frames the system is running on.

\begin{st}
3 class                  "SmallInteger"
3 class class            "SmallInteger class"
true class               "True"
nil class                "UndefinedObject"
[:x | x] class           "BlockClosure"
Object class class       "Metaclass"
\end{st}

\section{And the only verb is \texttt{send}}

Every computation is one object asking another to do something:

\begin{st}
3 + 4                    "send #+ to 3, with 4 as the argument"
anArray at: 1 put: 'x'   "send #at:put: to anArray"
Object subclass: #Point  "send #subclass: to a class -- classes are objects"
\end{st}

Even arithmetic. \ct{3 + 4} looks up \ct|#+| in \ct{SmallInteger} and runs
what it finds. The interpreter has fast paths for the common cases, but that
is an optimisation of a real message send, not a special case in the language:
you can override \ct|#+| in your own class and it works.

\subsection{What happens on a send}

\begin{enumerate}
\item The receiver's class is found.
\item Its method dictionary is searched for the selector.
\item If it is not there, the search continues in the superclass, and so on up
  to \ct{Object}, whose superclass is \ct{nil}.
\item If nothing implements it, \ct{doesNotUnderstand:} is sent instead, with
  the whole message reified as an object.
\end{enumerate}

That last step is worth dwelling on, because it is a language feature rather
than an error path. Chapter~\ref{ch:classes} uses it.

\section{\texttt{self} and \texttt{super}}

\ct{self} is the receiver. It is dynamically bound: in a method inherited from
\ct{Object}, \ct{self} is whatever the actual receiver is, not an
\ct{Object}.

\ct{super} is the \emph{same object}, but tells the lookup to start above the
class the currently-executing method was found in. It is not "my superclass's
instance"---there is only ever one object.

\begin{st}
printOn: aStream
    super printOn: aStream.          "let Object do its half"
    aStream nextPutAll: ' with '; print: self size
\end{st}

\begin{stwarn}
\ct{super} starts the search above the \emph{defining} class, not above the
receiver's class. Getting this wrong produces infinite recursion in exactly
one situation: writing \ct{self foo} in \ct{Foo>>foo} when you meant
\ct{super foo}.
\end{stwarn}

\section{Identity and equality}

Two different questions, two different selectors, and the difference matters
more in Smalltalk than in most languages because you can change what \ct{=}
means.

\begin{center}
\small
\begin{tabular}{@{}ll@{}}
\toprule
\ctp{==} & the same object. Never overridden.\\
\ctp{=}  & equal in value. Defined by the class, and often overridden.\\
\ctp{\~{}\~{}} & not the same object\\
\ctp{\~{}=} & not equal\\
\bottomrule
\end{tabular}
\end{center}

\begin{st}
'abc' = 'abc'            "true  -- same characters"
'abc' == 'abc'           "false -- two String objects"
#abc == #abc             "true  -- symbols are unique"
3 = 3.0                  "true"
3 == 3                   "true  -- SmallIntegers are their own identity"
\end{st}

The rule when you override \ct{=}: \textbf{override \ct{hash} too}, and make
sure that \ct{a = b} implies \ct{a hash = b hash}. A \ct{Set} or
\ct{Dictionary} that disagrees with you about hashing will lose your objects
in a way that is very hard to debug.

\section{\texttt{nil}}

\ct{nil} answers \ct{true} to \ct{isNil} and is what an uninitialised
instance variable holds. This system provides the modern conditionals:

\begin{st}
x ifNil: [0]                     "0 if x is nil, else x"
x ifNotNil: [:v | v printString]
x ifNil: ['none'] ifNotNil: [:v | v printString]
x isNil ifTrue: [...]            "the 1983 spelling; still fine"
\end{st}

\section{Printing and displaying}

Two protocols, and confusing them is a rite of passage.

\begin{center}
\small
\begin{tabular}{@{}lp{0.62\textwidth}@{}}
\toprule
\ctp{printString} & for a programmer: a form you could paste back into code.
  \ct{'abc' printString} is seven characters, including the quotes.\\[3pt]
\ctp{displayString} & for a user: the thing itself.
  \ct{'abc' displayString} is three characters.\\
\bottomrule
\end{tabular}
\end{center}

You customise both by overriding one method, \ct{printOn:}:

\begin{st}
printOn: aStream
    aStream nextPutAll: 'a Point ('; print: x; nextPutAll: ', '; print: y; nextPut: $)
\end{st}

\ct{printString} is defined in terms of \ct{printOn:}, so overriding the
stream version gets you the string version, and gets you correct behaviour
when your object is printed inside a collection.

\begin{stnote}[Some printStrings here are 1983's, not Pharo's]
\ct|#foo printString| answers \ct{'foo'}, not \ct|'#foo'|. \ct|#(1 2) printString| answers \ct{'(1 2 )'}, with a trailing space and no \ct|#|.
\ct{(3/4) printString} answers \ct{'(3/4)'} with parentheses. These are the
1983 library's own, and they are what the Browser and Inspector show you.
\end{stnote}

\section{Asking objects about themselves}

\begin{st}
3 class                          "SmallInteger"
3 isKindOf: Number               "true  -- this class or any superclass"
3 isMemberOf: SmallInteger       "true  -- exactly this class"
3 respondsTo: #printString       "true"
3 isNumber                       "true"
Object selectors size            "how many methods Object defines"
Integer includesSelector: #factorial
\end{st}

\ct{isKindOf:} is the one you usually want. Reaching for it often, though, is
a sign that you are writing a case statement where you should be writing
polymorphism.

\section{Polymorphism instead of conditionals}

This is the design lesson the language keeps trying to teach. Rather than
asking an object what it is and branching:

\begin{st}
"don't"
(shape isKindOf: Circle)
    ifTrue: [area := shape radius squared * Float pi]
    ifFalse: [(shape isKindOf: Square)
        ifTrue: [area := shape side squared]]
\end{st}

\ldots ask it to do the thing, and let each class answer for itself:

\begin{st}
"do"
area := shape area
\end{st}

with \ct{area} implemented in \ct{Circle} and in \ct{Square}. The whole
library is built this way, which is why \ct{do:} works on an \ct{Array}, an
\ct{OrderedCollection}, a \ct{Set}, an \ct{Interval} and a \ct{Stream}
without any of them knowing about the others.

\section{\texttt{doesNotUnderstand:}}

When lookup fails, the system sends \ct{doesNotUnderstand:} with a
\ct{Message} object holding the selector and arguments. The default
implementation raises \ct{MessageNotUnderstood}, which you can catch:

\begin{st}
[nil foo] on: MessageNotUnderstood do: [:e | e message selector]    "#foo"
\end{st}

Overriding it in your own class gives you proxies, loggers, and
domain-specific dispatch, and it is a supported technique rather than a hack.
Do it with restraint---a class that understands everything is a class no
Browser can help you with.
